% LNCS Template for Euro-Par Poster/Demo Paper
\documentclass[runningheads]{llncs}

% Required packages
\usepackage[T1]{fontenc}
\usepackage{graphicx}
\usepackage{amsmath,amssymb,amsfonts}
\usepackage{booktabs}
\usepackage{hyperref}

\begin{document}

\title{Sparse-Aware Tensor Parallelism for Efficient LLM Inference:\\
A Three-Level Solver Approach}

\author{Anonymous Author(s)}
\institute{Anonymous Institution}

\maketitle

\begin{abstract}
Tensor Parallelism (TP) is critical for distributed Large Language Model (LLM) inference. Existing approaches suffer from inter-GPU load imbalance, intra-GPU hardware underutilization, and excessive communication overhead. We present a sparse-aware TP framework with a three-level solver: (1) neuron partitioning based on co-activation patterns, (2) TensorCore/CUDACore assignment for hardware optimization, and (3) pipeline scheduling with sparse communication. Integration with DejaVu predictor and FlashFFN kernel achieves up to 2.3$\times$ speedup on OPT-175B inference while maintaining accuracy.

\keywords{Tensor Parallelism \and Large Language Models \and GPU Optimization}
\end{abstract}

\section{Introduction}

Large Language Model inference relies on Tensor Parallelism for multi-GPU distribution. Current approaches like Megatron-LM~\cite{megatron} and DeepSpeed~\cite{deepspeed} partition neurons uniformly, ignoring activation sparsity. The ``Lazy Neuron Phenomenon''~\cite{lazyneuron} shows over 90\% of neurons remain inactive, causing load imbalance and wasted computation.

Recent sparse prediction systems (DejaVu~\cite{dejavu}, PowerInfer~\cite{powerinfer}) enable accurate neuron activation forecasting. We leverage this to present a sparse-aware TP framework addressing three inefficiencies: (1) inter-GPU load imbalance via co-activation-aware partitioning, (2) intra-GPU hardware underutilization via TensorCore/CUDACore assignment, and (3) communication overhead via sparse communication primitives.

\section{Three-Level Solver Design}

\textbf{Level 1: Neuron Partitioning.} Given offline profiling mask $\mathbf{M}$, we construct co-activation matrix:
\begin{equation}
w_{uv} = \frac{1}{S}\sum_{i=1}^{S} M_{i,u} \cdot M_{i,v}
\end{equation}

Spectral clustering classifies neurons into Highly Co-activated Neuron Sets (HCNS) and Low-activation Neuron Sets (LCNS). Co-activated neurons reside on the same GPU to minimize cross-GPU communication, while ensuring balanced distribution across devices.

\textbf{Level 2: TC/CC Assignment.} We model execution time for TensorCore (TC) and CUDACore (CC):
\begin{equation}
T_{TC} = \frac{F^{TC} \cdot d \cdot s}{\text{BW}_{TC}} + \frac{2 \cdot F^{TC} \cdot d \cdot s}{\text{FLOPs}_{TC}}
\end{equation}
\begin{equation}
T_{CC} = \frac{F^{CC} \cdot d \cdot s_{active}}{\text{BW}_{CC}} + \frac{2 \cdot F^{CC} \cdot d \cdot s_{active}}{\text{FLOPs}_{CC}}
\end{equation}

We minimize $|T_{TC} - T_{CC}|$ via heuristic search, mapping HCNS to TensorCore and LCNS to CUDACore for optimal hardware utilization.

\textbf{Level 3: Pipeline \& Sparse Communication.} Input is divided into chunks for two-stage pipeline: computation of chunk $k+1$ overlaps with sparse AllReduce of chunk $k$. Sparse communication transmits only active (value, index) pairs, reducing volume from $O(s \cdot d)$ to $O(s_{active} \cdot d)$.

\section{Implementation}

\textbf{FlashFFN Kernel.} Fused kernel combining HCNS (TensorCore) and LCNS (CUDACore) computations with tiled processing and shared memory optimization.

\textbf{Integration with DejaVu.} Offline: profile activations and run three-level solver. Online: DejaVu predicts mask, FlashFFN executes with mask-aware computation, sparse communication aggregates results.

\section{Evaluation}

\textbf{Setup:} 8$\times$ NVIDIA A100 GPUs, OPT-175B and GPT-J-6B models. Baselines: Megatron-LM, DeepSpeed, DejaVu.

\begin{table}[h]
\caption{Inference latency (ms) and speedup}
\label{tab:results}
\centering
\begin{tabular}{lccc}
\toprule
\textbf{Method} & \textbf{OPT-175B} & \textbf{GPT-J-6B} & \textbf{Speedup} \\
\midrule
Megatron-LM & 125.4 & 18.3 & 1.0$\times$ \\
DejaVu & 89.2 & 14.1 & 1.4$\times$ \\
\textbf{Ours} & \textbf{54.7} & \textbf{9.2} & \textbf{2.3$\times$} \\
\bottomrule
\end{tabular}
\end{table}

\textbf{Ablation:} Level 1 contributes 1.4$\times$, Level 2 brings 1.8$\times$, Level 3 achieves 2.3$\times$ total speedup. Accuracy maintained within 1\% of dense baseline (perplexity 12.39 vs 12.34 on WikiText-2). Near-linear scaling up to 16 GPUs.

\textbf{Microbenchmarks:} Load variance reduced 45\%$\rightarrow$8\%, TC utilization 35\%$\rightarrow$72\%, CC utilization 15\%$\rightarrow$65\%, communication volume reduced 7.8$\times$.

\section{Conclusion}

Our three-level solver addresses key TP inefficiencies through co-activation-aware partitioning, heterogeneous hardware mapping, and sparse communication. Results demonstrate significant improvements while maintaining accuracy.

\bibliographystyle{splncs04}
\bibliography{references}

\end{document}
